\section{Coordinate and symmetry counterexamples}
\label{sec:symmetry}
Reflection is an additional restriction. The field $\psi_d=1+i_q^2$,
$\psi_q=i_q$ has d-even/q-odd parity but curl $-2i_q$. Thus symmetry alone
cannot establish a potential. For a declared even potential,
\begin{equation}
 \coenergy(i_d,i_q)=\coenergy(i_d,-i_q),\label{eq:energy-parity}
\end{equation}
its derivatives give the usual flux parities. This optional reflection is
not required for the integrability test.

A coherent change of orthogonal current coordinates is a contrasting case.
With $\current=\rotationmatrix(\alpha)\hat\current$, the chain rule yields
\begin{equation}
 \nabla_{\hat\current}\coenergy(\rotationmatrix(\alpha)\hat\current)
 =\rotationmatrix(-\alpha)\flux(\rotationmatrix(\alpha)\hat\current).
 \label{eq:rotated-potential}
\end{equation}
The Hessian undergoes symmetric congruence. Reciprocity survives while the
reflection axis rotates. Rotating only flux values while retaining old current
arguments is a different operation. These examples show what curl reveals
and what it cannot establish: a zero value tests conservation in the selected
conjugate coordinates, not alignment of a magnetic symmetry axis.
